% Section 1 — Conversion Overview
% TLDR for linguist reviewers. Proto-examples for pending specs,
% real examples for completed specs, technique comparison, status table,
% and related work on the colonial recorder perception problem.

\subsection{Conversion Examples and Methodologies}
\label{sec:examples}

This section presents concrete orthography-to-phoneme conversions across the
eight source languages. For sources where implementation is complete, real
conversion results are shown. For sources where only research has been
completed, \textbf{[PROTO]} examples are given --- these represent the current
best hypothesis grounded in the research evidence, with documented uncertainty.

\subsubsection*{Narragansett (Williams 1643) --- \texttt{[PROTO]}}

\begin{quote}
\textbf{Input:} \texttt{w\'etu} (house) \\
\textbf{Output:} /wi:tə w/ \\
\textbf{Method:} Comparative calibration via PA *wi:kiwa:mi, Massachusett
cognate \textit{wetu} (Eliot). The initial \textit{w\'e-} maps to /wi:/ via
Williams' circumflex-as-length convention; final \textit{-tu} maps to /tə w/
via PA reflex. \\
\textbf{Uncertainty:} Vowel length in the first syllable is inferred from the
circumflex accent, which Williams uses for primary stress --- whether this
marks phonemic length or stress prominence is debated.
\end{quote}

\begin{quote}
\textbf{Input:} \texttt{Aaunchem\'okaw} (Record 3834) \\
\textbf{Output:} /ã:čimo:hka:w/ \texttt{[PROTO]} \\
\textbf{Method:} FSM with preaspiration inference. Williams writes \textit{k}
but the only phonetic record (\textbackslash ph field) shows /hk/. The
preaspiration gap --- English-speaking recorders routinely fail to hear /h/
before stops --- is the single most important finding for this source. \\
\textbf{Uncertainty:} The nasalization of the initial vowel is inferred from
the \textit{Aaun-} spelling, which may represent /ã:/ or /aun/.
\end{quote}

\subsubsection*{Wampanoag (Trumbull/Eliot 1903) --- \texttt{[PROTO]}}

\begin{quote}
\textbf{Input:} \texttt{matta} (no, not) \\
\textbf{Output:} /mata/ \\
\textbf{Method:} Direct Eliot orthography mapping. Eliot's system is the most
regularized colonial orthography in the corpus. Confirmed via PA *mati-.
\end{quote}

\begin{quote}
\textbf{Input:} \texttt{achmōwonk} (devil) \\
\textbf{Output:} /čomu:wa\ng k/ \texttt{[PROTO]} \\
\textbf{Method:} 4-stage FSM pipeline: (1) normalize ō (Eliot's oo-ligature) to /u:/, (2) map
\textit{ach-} to /č/ via Eliot's c-before-a rule, (3) resolve the
\textit{-wonk} suffix via PA *-a:kani, (4) comparative fallback against
Pokanoket \textit{hobbamock} (Winslow) to check for dialect divergence. \\
\textbf{Uncertainty:} The ō symbol (Eliot's oo-ligature,
U+A74F) was historically misidentified as an infinity symbol. Its phonetic
value is long /u:/, but the precise quality (tense vs.\ lax) is debated.
\end{quote}

\subsubsection*{Mohegan-Pequot (Fielding 2013) --- \texttt{[PROTO]}}

\begin{quote}
\textbf{Input:} \texttt{wetu} (house) \\
\textbf{Output:} /wetu/ \\
\textbf{Method:} Direct 1:1 phoneme mapping. Fielding's modern practical
orthography is essentially phonemic. Confirmed via PA *wi:kiwa:mi and
cognate with Narragansett \textit{w\'etu} and Massachusett \textit{wetu}.
\end{quote}

\subsubsection*{Wampanoag (Wood 1634) --- \texttt{[PROTO]}}

\begin{quote}
\textbf{Input:} \texttt{abbona} (five) \\
\textbf{Output:} /apana/ \texttt{[PROTO]} \\
\textbf{Method:} Normalize \textit{b} to /p/ (Wood writes \textit{b} where
Eliot/Trumbull use \textit{p} --- direct evidence of voicing neutralization).
Segment the initial \textit{a-} as a demonstrative prefix. \\
\textbf{Uncertainty:} Wood's triple spelling for 'devil'
(\textit{abbomocho}/\textit{abomocho}/\textit{abamacho}) is the strongest
evidence in the corpus for colonial vowel instability. The medial vowel
could be /a/, /ə/, or /o/.
\end{quote}

\subsubsection*{Mahican (Edwards 1787) --- \texttt{[PROTO]}}

\begin{quote}
\textbf{Input:} \texttt{ghusooh} (eight) \\
\textbf{Output:} /xə sə w/ \texttt{[PROTO]} \\
\textbf{Method:} The \textit{gh} grapheme is the most unusual across all eight
sources. Candidates include /x/ (voiceless velar fricative), /h/, or /k/.
PA *nešwa:šika and Wampanoag cognate \textit{shwoshuk} suggest /x/
is most likely --- a Mahican phoneme not present in SNEA languages. \\
\textbf{Uncertainty:} The \textit{gh} value cannot be confirmed without
cross-linguistic cognate evidence from Munsee Delaware or other Mahican-adjacent
languages.
\end{quote}

\subsubsection*{Mohegan-Pequot (Prince-Speck 1904) --- \texttt{[PROTO]}}

\begin{quote}
\textbf{Input:} \texttt{bahkeder} \\
\textbf{Output:} /pahketə r/ \texttt{[PROTO]} \\
\textbf{Method:} German orthography normalizer: \textit{b-} initial = /p/
(German final devoicing rule applied to perception), \textit{-hk-} = /hk/
(preaspiration preserved --- Prince-Speck is one of the few sources that
writes preaspiration explicitly), \textit{-der} = /tə r/ (German
\textit{d} = /t/ in final position). \\
\textbf{Uncertainty:} The \textit{-er} suffix may represent /ə r/ or
a reduced syllable. Loanword detection needed --- \textasciitilde20--30\% of
the 446 entries are English loanwords.
\end{quote}

\subsection{Technique Comparison and Selection Rationale}
\label{sec:techniques}

The technique selection for each source is itself a research output, not a
predetermined input. Each source's unique characteristics --- corpus size,
orthographic regularity, ambiguity profile, available comparative evidence,
recorder linguistic background --- determine which approach is appropriate.
A one-size-fits-all approach would fail because the sources span 263 years
of colonial documentation with radically different orthographic systems.

Table~\ref{tab:techniques} compares the approaches considered across all
sources, with expected accuracy and rationale for selection or rejection.
The per-source rationale is documented in detail in Part~I.

\begin{table}[ht]
\centering
\caption{Technique comparison across SNEA sources}
\label{tab:techniques}
\small
\begin{tabular}{lp{3cm}p{2.5cm}p{4.5cm}}
\toprule
\textbf{Technique} & \textbf{Sources} & \textbf{Expected Accuracy} & \textbf{Rationale} \\
\midrule
Regex-based conversion & None & Low &
Rejected --- ambiguities are fundamentally context-dependent; regex cannot
distinguish /a/ from /ə/ from /a:/ in Williams' orthography without
comparative evidence. \\
\midrule
FSM with weighted transitions & Narragansett, Wampanoag, Mahican & Moderate &
Selected --- handles context-dependent ambiguity via state transitions;
weights encode probability from PA cognate frequency. \\
\midrule
1:1 phoneme mapping & Mohegan-Pequot (Fielding) & High &
Selected --- Fielding's modern practical orthography is already phonemic;
conversion is a direct lookup. \\
\midrule
ML classifier & Narragansett & Unknown &
Selected but blocked --- only 3 records have phonetic (\textbackslash ph)
fields; training data insufficient without hand-annotating \textasciitilde50
entries. \\
\midrule
Comparative calibration & All colonial sources & High for confirmed cognates &
Selected --- PA reconstructions and sister-language evidence (Massachusett,
Mohegan) disambiguate vowel quality and preaspiration. \\
\midrule
4-stage FSM pipeline & Wampanoag (Trumbull) & Moderate--High &
Selected --- preprocessing $\rightarrow$ phonemic conversion $\rightarrow$
comparative fallback $\rightarrow$ statistical verification. \\
\midrule
Loanword detection FSM & Mohegan-Pequot (Prince-Speck) & High &
Selected --- \textasciitilde20--30\% English loanwords must be identified
before phonemic conversion to avoid false cognate matches. \\
\midrule
German orthography normalizer & Mohegan-Pequot (Prince-Speck) & High &
Selected --- Prince's Neogrammarian training imposed German conventions
(\textit{tsh} for /ʃ/, \textit{z} for /ts/, Dehnungs-\textit{h}). \\
\midrule
Morphological segmentation & Mahican (Edwards) & Moderate &
Selected --- sentence-level entries require prefix/suffix parsing
(\textit{k-} 2nd person, \textit{n-} 1st person) before conversion. \\
\bottomrule
\end{tabular}
\end{table}

\subsection{Source Status and Expected KPIs}
\label{sec:status}

Table~\ref{tab:status} shows the current status of each source. Research is
complete for all eight sources. Implementation is pending for all sources.
Proto-KPI targets are research-grounded predictions based on the known
ambiguity profile of each source.

\begin{table}[ht]
\centering
\caption{Source status and proto-KPI targets}
\label{tab:status}
\small
\begin{tabular}{lcccc}
\toprule
\textbf{Source} & \textbf{Research} & \textbf{Implementation} &
\textbf{Expected Recall} & \textbf{Expected Accuracy} \\
\midrule
Narragansett (Williams 1643) & Complete & Pending & 85--90\% & 70--80\% \\
Wampanoag (Trumbull 1903) & Complete & Pending & 90--95\% & 80--90\% \\
Mohegan-Pequot (Fielding 2013) & Complete & Pending & 95--98\% & 90--95\% \\
Wampanoag multi-source & Complete & Pending & 85--90\% & 75--85\% \\
Wampanoag (Winslow 1624) & Complete & Pending & 70--80\% & 60--70\% \\
Wampanoag (Wood 1634) & Complete & Pending & 80--85\% & 65--75\% \\
Mohegan-Pequot (Prince-Speck 1904) & Complete & Pending & 85--90\% & 75--85\% \\
Mahican (Edwards 1787) & Complete & Pending & 80--85\% & 70--80\% \\
\bottomrule
\end{tabular}
\end{table}

\noindent\textbf{Note:} Expected accuracy is lower for sources with extreme
vowel ambiguity (Williams, Wood) or tiny corpora (Winslow). Mohegan-Pequot
(Fielding) is expected to have the highest accuracy because its modern
practical orthography is already phonemic. These targets will be updated as
implementation validates or refutes the research predictions.

\subsection{The Colonial Recorder Perception Problem}
\label{sec:related-work}

Every source in this study shares a fundamental challenge: English-speaking
recorders with different linguistic backgrounds hearing Algonquian phonemes
through their native phonological filter. This is not unique to SNEA --- it is
a well-documented phenomenon across colonial documentation of North American
languages.

\subsubsection*{Dr.\ Keith Cunningham --- Nanticoke}

Dr.\ Keith Cunningham (Linguist, Nanticoke Indian Tribe) completed his
doctoral dissertation \emph{A Phonological Analysis of Nanticoke With
Practical Applications for Language Revitalization} at Georgetown University
(2024). The dissertation analyzes three 18th-century Nanticoke word lists
(Heckewelder 1785) using the same methodology employed here: historical
phonology from colonial recorders, with the same challenges of multilingual
informants and orthographic complexity. His 2024 Algonquian Conference
presentation \emph{A Revised Phonological Analysis of the Heckewelder
Vocabulary of Nanticoke} directly parallels the analyses in this document.
Cunningham has also co-authored a Nanticoke language textbook, bridging
directly into practical revitalization --- the same goal as the SNEA
orthography-to-phoneme conversion pipeline.

Cunningham's work demonstrates that the colonial recorder perception problem
extends beyond SNEA into the broader Eastern Algonquian family (Nanticoke,
Maryland Algonquian). His methodology --- PA comparative calibration,
recorder-bias identification, cross-source verification --- is directly
transferable to the SNEA sources.

\subsubsection*{Dr.\ Craig Kopris --- Wyandot}

Dr.\ Craig Kopris (Independent scholar) has nearly three decades of experience
in language description and endangered language revitalization, with emphasis
on Iroquoian languages. His research note \emph{Les sons wyandots per\c{c}us
par des oreilles \'etrang\`eres} (Wyandot Sounds through Foreign Ears,
\emph{Recherches am\'erindiennes au Qu\'ebec}, 2014) is the exact framing of
the core problem addressed in this document: how European recorders
systematically misperceived and mis-transcribed indigenous phonemes. Kopris
demonstrates that knowing the sound system of the transcriber's first language
is the key to recovering the phonetic details of the target language.

His paper \emph{Wyandot Phonology: Recovering the Sound System of an Extinct
Language} (Proceedings of the Second Annual High Desert Linguistics Society
Conference, 1999) applies the same recovery methodology to Wyandot, a
Northern Iroquoian language unrelated to Algonquian. This demonstrates that
the foreign-ear perception problem is universal across colonial documentation
of North American languages, not limited to any single language family. His
Dictionary of Wyandot project (NEH award FN-255576-17) shows the long-term
application of these methods to language revitalization.

\subsubsection*{Relevance to This Project}

The work of Cunningham and Kopris establishes that the methodological
framework for addressing the colonial recorder perception problem is well
established. The SNEA orthography-to-phoneme project extends this framework
to eight sources across four languages (Narragansett, Wampanoag, Mohegan-Pequot,
Mahican) spanning 263 years of colonial documentation (1624--1904). The
techniques developed here --- FSM with weighted transitions, comparative
calibration via PA, colonial English phonology baselines, preaspiration
detection --- contribute to an ongoing methodological conversation about
recovering phonology from colonial sources across Eastern North America.
